\documentclass[conference]{IEEEtran}
\IEEEoverridecommandlockouts
\usepackage{cite}
\usepackage{amsmath,amssymb,amsfonts}
\usepackage{algorithmic}
\usepackage{graphicx}
\usepackage{textcomp}
\usepackage{xcolor}
\usepackage{hyperref}
\usepackage{booktabs}
\usepackage{listings}

\def\BibTeX{{\rm B\kern-.05em{\sc i\kern-.025em b}\kern-.08em
    T\kern-.1667em\lower.7ex\hbox{E}\kern-.125emX}}

\begin{document}

\title{An AI-Driven, Blockchain-Anchored Cyber Deception Framework for Adaptive Honeypot Defense}

\author{
\IEEEauthorblockN{Author Name}
\IEEEauthorblockA{
\textit{Department of Computer Science and Engineering} \\
\textit{Amrita Vishwa Vidyapeetham} \\
City, Country \\
email@example.com}
}

\maketitle

\begin{abstract}
Traditional honeypots -- Cowrie, Kippo, Dionaea, and similar low/medium-interaction
systems -- rely on static, deterministic response tables. This makes them
efficient at scale but fragile against fingerprinting: prior work has shown
that automated probing can distinguish such honeypots from real systems with
near-perfect accuracy in a single packet exchange \cite{vetterl2018bitter}. This
paper presents a multi-protocol cyber deception framework that replaces static
response tables with a large language model (LLM) response engine, producing
context-aware, session-consistent output for five decoy services (SSH, an
internal AI API, PostgreSQL, an Ethereum JSON-RPC node, and Jenkins CI/CD)
while retaining a fast static path for the $\approx$80\% of traffic that is
routine. Every interaction is intent-classified, SHA-256 hashed, and anchored
as a transaction on a local Ethereum chain, giving the resulting audit trail
tamper-evidence independent of the logging host's integrity. We evaluate the
system against a static-honeypot baseline along four axes -- sequential
latency/throughput, concurrent throughput, response diversity, and an
LLM-judge realism score -- across three successive evaluation runs, and
report the measured trade-offs explicitly rather than a single composite
metric. After an optimization pass (HTTP connection pooling applied across
the response engine and all five decoys, a tuned generation-length cap, and a
widened static fast-path), the system achieves a mean dynamic-path latency of
1.35~s (down from an initial 4.03~s, a $\approx$3$\times$ improvement)
against a 100-payload hand-written adversarial dataset, and a sequential
throughput of 0.74 requests/second, versus 3{,}333 requests/second for the
static fast-path. A third evaluation run replaced the hand-written dataset
with a 200-payload dataset built from real, publicly-documented attack
payloads (real SQLi, command-injection, XXE/Log4Shell/prompt-injection,
Ethereum JSON-RPC, and Jenkins-exploit strings; see
Section~\ref{sec:realworldrun}), which increased mean dynamic latency to
2.06~s -- a genuine effect of longer, more structurally varied real payloads,
not a regression -- and surfaced a real weakness invisible to the
hand-written set: response diversity on the API/prompt-injection decoy
dropped from 1.00 to 0.70 against real XXE/Log4Shell payload shapes. This
same run also uncovered and fixed a benchmark-harness bug that had silently
measured Redis cache-hit latency (8.2~ms, reported as ``800 requests/second'')
instead of genuine concurrent LLM throughput; the corrected
concurrent-throughput figure is \textbf{1.62 requests/second} at a mean
per-request latency of \textbf{3.09~s} under 10 simultaneous sessions.
Response diversity otherwise remained high (0.80--1.00 across four of five
services) versus near-zero for a static baseline by construction, and the
mean LLM-judged realism score was stable across all three runs
(4.00--4.03/5). We report each run's numbers, the harness bug, and its
correction explicitly rather than presenting only the most favorable
figures, and discuss where the resulting trade-offs are acceptable
(human-operated intrusion, not internet-wide scanning) and the limitations
of automated realism scoring as a proxy for human deception studies.
\end{abstract}

\begin{IEEEkeywords}
honeypot, cyber deception, large language models, blockchain, digital forensics, intrusion detection, adaptive defense
\end{IEEEkeywords}

\section{Introduction}

Honeypots are decoy systems deployed to attract, detect, and study
unauthorized access attempts without exposing production infrastructure.
Since Cowrie's introduction as an SSH/Telnet honeypot, low- and
medium-interaction honeypots have become a standard component of threat
intelligence pipelines: they are cheap to deploy, safe to run, and generate
large volumes of attacker telemetry. Their core weakness is also their core
design choice -- a finite table of canned responses. An attacker, or more
commonly an automated scanner, can send a handful of atypical commands and
observe whether the responses are internally consistent with a real system.
Vetterl and Clayton demonstrated Internet-scale fingerprinting of exactly this
class of honeypot with a single distinguishing packet and near-zero error rate
\cite{vetterl2018bitter}. Once a honeypot is identified, it stops collecting
useful intelligence on capable adversaries, who simply move on or feed it
false telemetry.

Two independent trends motivate a different approach. First, instruction-tuned
LLMs can generate plausible, contextually consistent text for narrow,
well-specified personas (a Linux shell, a SQL error, a JSON-RPC response) far
more cheaply than hand-authoring a comprehensive fake filesystem or database.
Recent work -- shelLM \cite{sladic2024shellm}, HoneyGPT \cite{wang2024honeygpt},
and related 2023--2025 studies \cite{mckee2023chatbots,llmhoneypot2024} -- has
shown that LLM-generated shell sessions can hold up under sustained human
interaction significantly better than static emulation. Second,
blockchain-anchored hashing offers a low-cost way to make forensic logs
tamper-evident: instead of trusting the honeypot host's own log integrity
(which an attacker who compromises the host could edit), a hash of each log
entry is committed to an independent, append-only ledger, so any later
modification is detectable by re-hashing and comparing.

This paper combines both trends into a single deployable framework and
contributes:

\begin{enumerate}
\item A \textbf{multi-protocol deception architecture} covering five distinct
attack surfaces (SSH, an internal AI API, SQL, Ethereum RPC, and CI/CD) behind
a single logging and classification pipeline -- broader coverage than the
SSH-only focus of most LLM-honeypot literature to date.
\item A \textbf{three-tier response strategy} (static dictionary
$\rightarrow$ cached LLM response $\rightarrow$ live LLM call) that keeps
routine traffic at millisecond latency while reserving expensive generation
for genuinely novel input.
\item \textbf{Blockchain-anchored log integrity} applied specifically to
honeypot telemetry, giving each log entry a verifiable, tamper-evident hash
independent of the collecting host.
\item An \textbf{evaluation methodology} that reports latency, response
diversity, and an LLM-judge realism score as separate, reproducible
measurements -- explicitly avoiding a single unexplained accuracy figure --
along with a discussion of what an automated judge can and cannot substitute
for.
\end{enumerate}

\section{Related Work}

\subsection{Static Honeypots and Fingerprinting}

Cowrie, Kippo, Dionaea, and T-Pot remain widely deployed low/medium-interaction
honeypots, valued for low resource cost and safety. Their static response
model is precisely what fingerprinting techniques exploit: Vetterl and Clayton
\cite{vetterl2018bitter} achieved an equal error rate of 0.0183 detecting such
honeypots at Internet scale using single-packet probes; a later multistage
active-probing framework \cite{gotcha2023} generalized this to a broader class
of honeypots. A 2026 survey of fingerprinting and evasion techniques
\cite{honeypotsurvey2026} confirms this remains an open problem rather than
one solved by incremental patching of response tables.

\subsection{LLM-Backed Deception}

Sladi\'{c} et al.'s shelLM \cite{sladic2024shellm} used an LLM to generate an
entire SSH shell session and reported a 0.90 true-negative rate -- human
evaluators judged the sessions to be a real shell 90\% of the time. HoneyGPT
\cite{wang2024honeygpt} extended this with chain-of-thought prompting for more
stateful sessions and reported deeper attacker engagement in field deployment.
McKee and Noever's early exploration \cite{mckee2023chatbots} and the broader
``LLM Honeypot'' study \cite{llmhoneypot2024} corroborate that LLM-driven
responses measurably outperform static emulation on interactivity and
perceived realism. This body of work is concentrated on SSH/terminal
honeypots; our framework applies the same idea across five protocol types
with a shared backend, which to our knowledge has not been evaluated together
in prior work.

\subsection{Blockchain for Log Integrity and Forensic Chain of Custody}

Hash-then-anchor patterns for tamper-evident logging are established in
digital forensics: a 2025 IEEE study on blockchain-based evidence tracking
\cite{blockchainevidence2025} and a Springer paper on tamper-proof cloud VM
investigation models \cite{cloudvmforensics2022} both apply hashing plus
blockchain anchoring to forensic evidence generally. A PeerJ Computer Science
paper on blockchain-enforced immutable auditability for threat detection
\cite{peerjthreat2025} applies a closely related pattern to security telemetry
specifically. We did not find prior work applying this pattern to honeypot
logs specifically -- our contribution is the application, not the underlying
cryptographic technique.

\subsection{Deception Frameworks and Honeytokens}

MITRE Engage \cite{mitreengage2022} provides a structured vocabulary for
adversary engagement and deception mapped to ATT\&CK, which we use loosely to
frame our five decoys as covering distinct engagement opportunities
(credential harvesting, prompt injection, data exfiltration bait,
transaction-signing bait, and supply-chain bait respectively). Honeytokens --
individually identifiable fake credentials seeded to detect misuse -- have a
long history dating to HoneyGen \cite{bercovitch2011honeygen}; our seeded fake
API keys and database records in \texttt{vault-sql} are a direct application
of this idea.

\subsection{Realism and Detectability Benchmarking}

Standardized telemetry metrics for honeypot realism (dwell time, interaction
depth, behavioral diversity) have been proposed recently
\cite{ijnsa2026telemetry}; we adopt ``response diversity'' as one such metric
in our own evaluation. We could not verify several secondary-source
dwell-time statistics circulating for RL/LLM-driven deception systems and
deliberately do not cite unverified numbers here -- see Section~\ref{sec:limitations}.

\subsection{Direct Comparison to Cited Systems}

Table~\ref{tab:comparison} places this framework against the honeypot
systems and papers cited above. ``Interaction level'' follows the standard
low/medium/high taxonomy (low = fixed banner/response, medium = scripted
protocol emulation with some state, high = a real or near-real backend).
Entries are included only where we could verify the claim from the system's
own documentation or the cited paper; unverified secondary-source figures
are marked as such rather than presented as fact.

\begin{table*}[htbp]
\caption{Comparison against cited honeypot systems.}
\label{tab:comparison}
\centering
\begin{tabular}{@{}p{2.1cm}p{2.6cm}p{4.6cm}p{2.6cm}p{2.6cm}@{}}
\toprule
\textbf{System} & \textbf{Interaction level} & \textbf{AI/LLM-generated responses} & \textbf{Protocols / services} & \textbf{Tamper-evident logging} \\
\midrule
Cowrie \cite{cowrie} & Medium (emulated shell) or high (proxy mode) & No (static/scripted); an experimental LLM fallback mode exists in recent releases but is a single-tier call, not a static$\to$cache$\to$LLM pipeline & SSH, Telnet & No -- plain JSON/text logs \\
T-Pot \cite{tpot} & Mixed -- orchestrates $\approx$20 independent honeypot daemons (incl. Cowrie, Dionaea) behind an ELK dashboard & No -- aggregation/visualization layer only, no response generation of its own & Broad (SSH, SMB, web, RDP, ICS/SCADA, etc.), via bundled daemons & No \\
Dionaea \cite{dionaea} & Low--medium (scripted protocol emulation for payload capture) & No & SMB, MySQL wire protocol, SIP, FTP, TFTP & No \\
shelLM \cite{sladic2024shellm} & High for SSH sessions specifically & Yes -- full LLM-generated shell session & SSH only & Not addressed \\
HoneyGPT \cite{wang2024honeygpt} & High for terminal sessions & Yes -- GPT-3.5/4 via a Cowrie-derived proxy; authors report an internally-defined $\approx$92\% attack-success and $\approx$96\% deception-accuracy rate (self-reported, not independently reproduced -- treat as a single-study figure) & SSH/terminal only & Not addressed \\
This framework & High across all five decoys & Yes -- 3-tier (static dict $\to$ Redis cache $\to$ LLM), session-history-aware, applied uniformly across 5 distinct protocol types & SSH, internal AI API, PostgreSQL, Ethereum JSON-RPC, Jenkins CI/CD & Yes -- SHA-256 per log entry, anchored to an Ethereum chain \\
\bottomrule
\end{tabular}
\end{table*}

Two gaps in the cited literature motivate this framework's design choices
rather than being incidental: first, LLM-backed deception work to date
(\cite{sladic2024shellm}, \cite{wang2024honeygpt}, \cite{mckee2023chatbots},
\cite{llmhoneypot2024}) is concentrated on a single protocol (SSH/terminal);
we are not aware of a prior system applying an LLM response engine uniformly
across five heterogeneous protocol types (shell, REST/chat API, SQL,
JSON-RPC, CI/CD webhook) behind one shared backend. Second, a
systematization-of-knowledge survey covering both honeypots and LLMs
\cite{sokhoneypotllm2025} does not discuss blockchain or tamper-evident
logging as a honeypot design axis at all, and we found no honeypot-specific
paper combining hash-anchoring with decoy logging -- the closest related
work is general (non-honeypot) tamper-evident audit-log research
\cite{blockchainevidence2025,cloudvmforensics2022,peerjthreat2025}. We
report this as an absence in the literature we searched, not as a claim that
no such system exists.

\section{System Architecture}

The framework has six components: a network-facing proxy, five decoy
services, a shared AI response engine, a central backend, a blockchain
ledger, and a live dashboard. Only the proxy and dashboard are exposed to the
network; all other components communicate over an internal Docker network,
so a compromised decoy cannot directly reach the logging pipeline or the
blockchain client.

\begin{figure}[htbp]
\centering
\fbox{\begin{minipage}{0.46\textwidth}
\footnotesize
\begin{verbatim}
INTERNET / ATTACKER
        |
        v
  +-------------+     logs every connection
  |    PROXY    |----------------------------+
  | (port guard)|                            |
  +-------------+                            v
        |                            +---------------+
  routes to 5 decoys:                |    BACKEND    |
        |                            |   (FastAPI)   |
  +-----+----------------------+     | classify/hash |
  | shadow-shell  (SSH)        |     +------+--------+
  | oracle-ai     (AI API)     |            |
  | vault-sql     (PostgreSQL) |     +------+---+-----------+
  | ledger-trap   (Ethereum)   |     |          |           |
  | forge-ci      (Jenkins)    |  REDIS     GANACHE     DASHBOARD
  +-----+----------------------+ (store) (blockchain)  (React/WS)
        |
        v
  AI ENGINE (OpenRouter-backed)
\end{verbatim}
\end{minipage}}
\caption{System architecture. Only the proxy and dashboard are
network-exposed; all other components communicate over an internal network.}
\label{fig:architecture}
\end{figure}

\subsection{Proxy and Network Isolation}

\texttt{proxy/proxy.py} is an \texttt{asyncio} TCP server listening
simultaneously on five public ports (2222, 8080, 5433, 8545, 9090). On each
new connection it asynchronously logs the source IP and target port to the
backend, then transparently pipes bytes between the attacker and the
corresponding internal decoy. This separates connection-level telemetry
(useful even for port scanners that never send a payload) from
application-level decoy logic, and ensures every decoy's internal listening
port is never directly reachable from outside the Docker network.

\subsection{Decoy Services}

Five decoys cover distinct attack surfaces, each impersonating a specific
real system and each logging every interaction to the backend before (or
while) generating a response:

\begin{itemize}
\item \textbf{shadow-shell} (SSH, port 2222) -- a real \texttt{asyncssh}
server that accepts any username/password combination, logs every login
attempt, and presents a shell prompt for \texttt{dbadmin@corp-prod-db-01}.
Directory navigation (\texttt{cd}) is handled locally for speed and
consistency; all other commands are routed through the AI engine.
\item \textbf{oracle-ai} (internal AI API, port 8080) -- a FastAPI service
exposing \texttt{/api/chat}, \texttt{/api/predict}, and \texttt{/api/models}
endpoints for a fictitious internal model (\texttt{gpt-corp-v2}).
Prompt-injection attempts (keywords like ``jailbreak,'' ``ignore previous
instructions'') are logged and answered with a scripted ``bypass accepted''
response, keeping the attacker engaged and their injection technique fully
captured.
\item \textbf{vault-sql} (PostgreSQL decoy, port 5433) -- a FastAPI query
endpoint plus a real, separately-seeded PostgreSQL instance containing
synthetic user records, financial records, and API keys formatted to
resemble real OpenAI, GitHub, and AWS credentials, used as honeytokens.
\item \textbf{ledger-trap} (Ethereum JSON-RPC, port 8545) -- responds to
standard Web3 RPC calls (\texttt{eth\_getBalance}, \texttt{eth\_accounts},
etc.) with statically-defined plausible values for common calls and routes
uncommon RPC methods through the AI engine.
\item \textbf{forge-ci} (Jenkins CI/CD, port 9090) -- exposes Jenkins-style
endpoints (\texttt{/api/json}, \texttt{/job/build/api/json},
\texttt{/git/notifyCommit}) that fake a working build/deploy pipeline,
targeting supply-chain attackers.
\end{itemize}

\subsection{Three-Tier AI Response Engine}

\texttt{ai-engine/engine.py} answers each request in one of three tiers, in
order, so that the expensive LLM call is reserved for genuinely novel input:

\begin{enumerate}
\item \textbf{Static dictionary.} Roughly a dozen common SSH commands
(\texttt{ls}, \texttt{whoami}, \texttt{env}, \texttt{ps aux},
\texttt{cat /etc/passwd}, etc.) and structural patterns (\texttt{cd },
\texttt{sudo }, \texttt{echo }) are answered instantly from a fixed table --
this is the same response model a traditional static honeypot uses for
everything, and remains the fast path here for routine traffic.
\item \textbf{Redis cache.} Unknown commands are hashed (including up to two
prior history entries, so repeated attacks in a similar context reuse a
cached answer) and checked against a Redis cache with a multi-hour TTL
before any LLM call is made.
\item \textbf{Hosted LLM call.} A cache miss is sent to a hosted
chat-completion model via the OpenRouter API, with a service-specific system
prompt instructing the model to output only the raw response a real system
would produce -- no explanation, no markdown, no acknowledgment that it is a
language model -- and, for the SSH service, session history is included so
the shell exhibits state consistency (e.g., a file created in one command
persists in a later \texttt{ls}).
\end{enumerate}

This mirrors the three-tier design validated by prior LLM-honeypot work
(e.g., shelLM, HoneyGPT) while extending it across five protocol types
sharing one engine and one cache.

\subsection{Intent Classification}

\texttt{backend/main.py::classify\_intent} applies keyword matching to every
logged payload, assigning one of eight categories: SQL injection, privilege
escalation, reverse shell, cryptomining, prompt injection, supply-chain,
reconnaissance, or unknown. This is intentionally simple -- a rule-based
baseline rather than a learned classifier -- and is reported as a known
limitation (Section~\ref{sec:limitations}) rather than a contribution in its
own right; it exists to make the live dashboard immediately useful, not as a
novel detection technique.

\subsection{Blockchain-Anchored Logging}

Every log entry received at \texttt{POST /api/log} is: (1) timestamped and
intent-classified; (2) serialized with sorted keys and SHA-256 hashed; (3)
anchored by sending a zero-value self-transaction on a local Ethereum chain
(Ganache) with the hash as the transaction's \texttt{data} field; (4) stored
in a bounded Redis list (last 10{,}000 entries) alongside the hash and
anchoring status; and (5) broadcast over WebSocket to connected dashboard
clients. Because the hash is committed to an independently operated ledger,
any later modification of the stored log is detectable by re-hashing it and
comparing against the anchored value -- Section~\ref{sec:blockchaintest}
reports a direct test of this property.

\subsection{Real-Time Dashboard}

A React single-page app (\texttt{dashboard/src/App.jsx}) polls
\texttt{/api/stats} every five seconds and subscribes to \texttt{/ws/live}
for a push feed, showing aggregate stats (total attacks, unique IPs, intent
breakdown, top attacker IPs) and a scrolling live feed of the last 200
events, each annotated with its blockchain-anchoring status.

\section{Evaluation Methodology}

We evaluate against a \textbf{static-honeypot baseline}, operationalized as
the system's own static fast-path (a deterministic dictionary lookup,
functionally equivalent to Cowrie's canned-response model) versus its
LLM-backed dynamic path, using \texttt{evaluate.py} against the live system.

\subsection{Dataset}

Results in Sections~\ref{sec:baseline}--\ref{sec:blockchaintestresults} use a
100-payload hand-curated adversary dataset
(\texttt{apt\_attack\_dataset.json}), 20 payloads per decoy service,
generated by an earlier version of \texttt{generate\_dataset.py} from
hand-written attack patterns (reverse shells, SQL injection, RPC abuse,
prompt injection, CI/CD supply-chain payloads) each tagged with a random
nonce to defeat caching during evaluation. Section~\ref{sec:realworldrun}
reports a subsequent run using a 200-payload dataset (40 per service)
rebuilt from real, publicly-documented attack payload corpora rather than
hand-written strings -- see Section~\ref{sec:realworldrun} and the project
README for the full source list and rationale for the switch.

\subsection{Metrics}

\begin{enumerate}
\item \textbf{Latency / throughput.} 30 trials of the static fast-path and
20 trials of the LLM-backed path; mean, median, and standard deviation
reported rather than a single sample, addressing a methodological gap in the
project's original benchmark (which used a single dynamic-path measurement
and reported it as if representative).
\item \textbf{Response diversity.} Across the full 100-item dataset, the
fraction of distinct outputs per service is used as a proxy for resistance
to signature-based fingerprinting -- a static honeypot necessarily scores at
or near the number of distinct hardcoded templates, however many inputs it
receives.
\item \textbf{LLM-judge realism score.} An independent model
(\texttt{meta-llama/llama-3.1-70b-instruct}, distinct from the
\texttt{meta-llama/llama-3.1-8b-instruct} response-generation model to
reduce self-evaluation bias) rates each response 1--5 against a fixed
rubric, on a sample of 30 payloads.
\end{enumerate}

\subsection{Limitations of the Realism Metric}

This is an \textbf{automated proxy metric}, not a human-subject study. LLM
judges are known to have their own biases (e.g., favoring longer or more
verbose responses) and may not perfectly track what a human attacker would
find convincing. shelLM's 0.90 true-negative rate \cite{sladic2024shellm}
was obtained with human evaluators; we report our LLM-judge score as a
weaker but far cheaper and more reproducible substitute, and explicitly
recommend a human evaluation as follow-up work rather than presenting this
score as equivalent evidence.

\subsection{Blockchain Integrity Test}
\label{sec:blockchaintest}

We additionally demonstrate tamper-evidence directly
(\texttt{verify\_integrity.py}): a real log entry is submitted through the
live pipeline and anchored on-chain; it is then fetched back, re-hashed, and
confirmed to match its recorded hash; a modified copy of the same entry is
then re-hashed and confirmed to \emph{not} match, demonstrating that the
hash anchored at logging time would catch a post-hoc modification.

\section{Results}

All figures below are measured directly against the live local stack
(\texttt{evaluate.py}, \texttt{verify\_integrity.py}) using
\texttt{meta-llama/llama-3.1-8b-instruct} as the response-generation model
and \texttt{meta-llama/llama-3.1-70b-instruct} as the judge; full raw output
is available in \texttt{evaluation\_results.json}.

\subsection{Initial Baseline}
\label{sec:baseline}

The first evaluation run, before any latency optimization, established the
baseline in Table~\ref{tab:baseline}.

\begin{table}[htbp]
\caption{Initial Baseline: Static vs. LLM-Backed Path}
\label{tab:baseline}
\centering
\begin{tabular}{@{}lll@{}}
\toprule
\textbf{Metric} & \textbf{Static} & \textbf{Dynamic} \\
\midrule
Mean latency & 0.4 ms & 4{,}028 ms (4.03 s) \\
Median latency & 0.4 ms & 3{,}529 ms (3.53 s) \\
Stdev & 0.1 ms & 2{,}380 ms (2.38 s) \\
Min / Max latency & 0.3 / 0.7 ms & 1{,}037 / 8{,}807 ms \\
Throughput (req/s) & 2{,}500 & 0.25 \\
Diversity ratio (mean) & $\approx$0 & 0.99 (0.95--1.00) \\
Avg. response length (words) & 0 & 16.2 \\
\bottomrule
\end{tabular}
\end{table}

LLM-judge realism score: mean 4.38/5 ($\sigma = 1.28$, $n=24$ of 30 sampled
responses successfully parsed -- six judge calls returned unparseable output
during this run and are excluded rather than assumed either way).

\subsection{Optimization Pass}
\label{sec:optimization}

Profiling the baseline identified two concrete, fixable inefficiencies in
\texttt{ai-engine/engine.py}: (1) a new \texttt{httpx.AsyncClient} -- and
therefore a fresh TCP+TLS handshake -- was opened on \emph{every single
request} to OpenRouter instead of reusing a connection pool; (2)
\texttt{max\_tokens=150} was generous relative to what these responses
need, and generation length is the dominant cost in a chat-completion
request's latency. A third, lower-effort change widened the static SSH
dictionary with 12 additional common recon commands so more traffic never
reaches the LLM path at all.

Before applying the token-length change, a direct measurement (3 trials,
identical prompt) confirmed generation length as the dominant latency
factor: \texttt{max\_tokens=150} averaged $\approx$5.3~s,
\texttt{max\_tokens=80} $\approx$3.4~s, \texttt{max\_tokens=40}
$\approx$2.0~s. 60 was selected as a middle ground. A separate measurement
also tested whether a smaller model would help: comparing
\texttt{llama-3.1-8b-instruct} against the smaller
\texttt{llama-3.2-3b-instruct} and \texttt{llama-3.2-1b-instruct} (3 trials
each, identical prompt) found the 8B model was consistently the
\emph{fastest} of the three (1.3--1.9~s) and the 1B model the
\emph{slowest} (1.8--2.5~s) -- provider/routing behavior on a hosted API
dominates over raw parameter count, so a model downgrade was not applied.

A second, independent evaluation run (fresh Redis cache, same methodology)
after these changes measured the results in Table~\ref{tab:optimization}.

\begin{table}[htbp]
\caption{Latency/Throughput: Before vs. After Optimization}
\label{tab:optimization}
\centering
\begin{tabular}{@{}llll@{}}
\toprule
\textbf{Metric} & \textbf{Before} & \textbf{After} & \textbf{Change} \\
\midrule
Dynamic mean latency & 4.03 s & 1.35 s & $\approx$3$\times$ \\
Dynamic median latency & 3.53 s & 1.25 s & $\approx$3$\times$ \\
Dynamic stdev & 2.38 s & 0.63 s & tighter \\
Dynamic throughput (seq.) & 0.25 rps & 0.74 rps & $\approx$3$\times$ \\
Static throughput & 2{,}500 rps & 3{,}333 rps & +33\% \\
LLM-judge realism & 4.38/5 (24/30) & 4.00/5 (30/30) & $-$0.38 \\
\bottomrule
\end{tabular}
\end{table}

A new metric -- \textbf{concurrent throughput} -- was added for this second
run: firing 10 requests simultaneously (rather than one at a time). The
first measurement of this metric reported 800.3 requests/second at a mean
per-request latency of 8.2~ms; this figure was later discovered to be a
benchmark-harness bug, not a genuine measurement (\texttt{evaluate.py}'s
concurrent-throughput step reused payloads already sent once earlier in the
same run, so they were guaranteed Redis cache hits -- the 8.2~ms was cache
round-trip latency, not LLM latency). After fixing the harness to force a
fresh cache miss per concurrent request (Section~\ref{sec:realworldrun}),
the corrected figure is \textbf{1.62 requests/second} at a mean per-request
latency of \textbf{3.09~s} under 10 concurrent sessions. This remains the
more relevant number for real deployment capacity than the sequential
figures, since concurrent requests overlap their network round-trip time --
but it is three orders of magnitude lower than first reported, and the
earlier 800~req/s figure should not be cited.

The realism-score drop (4.38 $\rightarrow$ 4.00) is honestly reported as a
real cost of the shorter generation cap, though it is confounded with a
data-completeness improvement: the ``before'' score is computed over only
24 of 30 samples (6 failed to parse, due to a transient OpenRouter error
unrelated to the engine changes -- see Section~\ref{sec:limitations}),
while the ``after'' score reflects all 30. Judge reasons for the
lower-scored ``after'' responses cite missing detail (``lacks specific
details,'' ``missing expected Jenkins-specific details'') consistent with
the shorter token cap, not a change in the model's persona-consistency.
Response diversity (0.95--1.00) and the blockchain integrity property were
unaffected by these changes.

\subsection{Extending the Fix to All Five Decoys}

The ai-engine fix above addressed one hop in the request path (ai-engine
$\rightarrow$ OpenRouter). An identical anti-pattern existed one hop
earlier, in every one of the five decoys: each decoy's logging call (every
attacker interaction posted to the backend) and its ai-engine call opened a
new \texttt{aiohttp.ClientSession} -- a fresh TCP connection -- on every
single request rather than reusing one. This was fixed identically across
all five decoys (shadow-shell, oracle-ai, vault-sql, ledger-trap,
forge-ci), each now creating one pooled session at process startup instead
of per-request.

To isolate the decoy layer's own contribution from the ai-engine's, the
before/after code was checked out via git and run side-by-side on
different ports, against static/deterministic payloads so ai-engine
latency does not dominate the measurement (Table~\ref{tab:decoys}).

\begin{table}[htbp]
\caption{Decoy-Layer Connection Pooling: Before vs. After}
\label{tab:decoys}
\centering
\begin{tabular}{@{}llll@{}}
\toprule
\textbf{Decoy} & \textbf{Before (mean)} & \textbf{After (mean)} & \textbf{$\Delta$} \\
\midrule
oracle-ai & 6.15 ms ($\sigma$=3.89) & 4.28 ms ($\sigma$=0.70) & $-$30.4\% \\
vault-sql & 0.69 ms ($\sigma$=1.94) & 0.41 ms ($\sigma$=0.44) & $-$41.4\% \\
ledger-trap & 5.96 ms ($\sigma$=2.43) & 4.20 ms ($\sigma$=0.75) & $-$29.5\% \\
forge-ci & 6.95 ms ($\sigma$=3.13) & 6.71 ms ($\sigma$=3.02) & $-$3.4\% \\
shadow-shell & 7.91 ms ($\sigma$=9.94) & 4.64 ms ($\sigma$=6.55) & $-$41.3\% \\
\bottomrule
\end{tabular}
\end{table}

The effect is smaller in absolute terms than the ai-engine fix -- localhost
TCP handshakes are cheap relative to a cross-internet TLS handshake to
OpenRouter -- but real and consistent across four of five decoys, and it
substantially tightens latency variance (standard deviation drops by
roughly half to two-thirds in every case): a connect-per-request pattern
is more exposed to occasional slow handshakes than a warm, reused
connection. forge-ci shows negligible change because its measured endpoint
always routes through the ai-engine regardless of decoy-layer pooling, so
its latency floor is set by the ai-engine's own cache lookup, not the
decoy's connection handling.

\subsection{Blockchain Integrity Test}
\label{sec:blockchaintestresults}

A log entry with hash \texttt{356b7496...33a9a1} (truncated) was
submitted, anchored on-chain (\texttt{blockchain\_anchored: true}), and
fetched back; re-hashing the unmodified entry reproduced the identical
hash. Re-hashing a copy with \texttt{raw\_data} altered from
``\texttt{sudo cat /etc/shadow}'' to ``\texttt{ls -la}'' produced a
different hash (\texttt{7530fb00...0d9d394}), confirming the
tamper-evidence property: any change to a logged entry after anchoring is
detectable by recomputation, without needing to trust the host storing the
log. This test is unaffected by the optimization pass in
Section~\ref{sec:optimization}.

\subsection{Third Run: Real-World Dataset and a Benchmark-Harness Correction}
\label{sec:realworldrun}

The adversarial dataset used in
Sections~\ref{sec:baseline}--\ref{sec:blockchaintestresults}
(\texttt{apt\_attack\_dataset.json}) was originally hand-written:
representative SQLi, command-injection, and RPC-abuse strings authored to
\emph{look like} real attack traffic. To ground the evaluation in verifiably
real attacker input rather than plausible-looking synthetic strings, the
dataset was rebuilt to pull real, publicly-documented payloads from named
sources: real command-injection strings for \texttt{ssh}
(\texttt{payload-box/command-injection-payload-list}), real SQL injection
strings for \texttt{sql} (\texttt{payload-box/sql-injection-payload-list}),
real prompt-injection/SSRF/XXE/Log4Shell strings for \texttt{api}
(\texttt{swisskyrepo/PayloadsAllTheThings}), the documented Ethereum
JSON-RPC methods an exposed node's attack surface actually consists of for
\texttt{web3}, and documented Jenkins script-console/path-traversal
techniques for \texttt{cicd}. The dataset grew from 100 to 200 payloads
(40/service; see the project README's ``Dataset provenance'' section for
the full source list). \texttt{evaluate.py} was re-run in full against the
live stack with this dataset.

Two results are worth reporting on their own merits, independent of the
earlier optimization narrative:

\textbf{Response diversity on the \texttt{api} service dropped from 1.00 to
0.70} -- the clearest signal this re-run produced that the hand-written
dataset had missed. Inspecting the judge's per-sample reasoning
(\texttt{evaluation\_results.json}) shows the model frequently collapses to
a generic ``invalid request'' response when the payload is a real XXE
(\texttt{<!ENTITY ...>}) or Log4Shell (\texttt{\$\{jndi:...\}}) string,
rather than emulating a plausible vulnerable-parser response to that exact
payload -- e.g. \emph{``Generic error message without any specific details
or XML parsing error indication.''} The hand-written \texttt{api} payloads
were simple JSON chat-style prompts the model already handled fluently, so
this weakness was invisible until real payload shapes were used. It is now
a concrete, evidence-backed follow-up: the \texttt{oracle-ai} decoy's
system prompt likely needs XML/JNDI-aware few-shot examples.

\textbf{A genuine bug in the evaluation harness itself was found and fixed
during this re-run.} \texttt{bench\_concurrent\_throughput()} had reused
\texttt{dataset[:CONCURRENCY\_LEVEL]} -- the same 10 payloads already sent
once earlier in the same run via \texttt{diversity\_pass()}, which
exercises the entire dataset. Because the ai-engine's Redis cache key is
\texttt{md5(command + history)} (\texttt{ai-engine/engine.py}), these were
guaranteed cache hits by the time the concurrent-throughput step ran,
despite the function's docstring asserting the payloads were cache-proof.
They were not: the dataset's nonces are generated once at dataset-build
time, not refreshed per benchmark call. This means \textbf{the 800.3
requests/second concurrent-throughput figure reported in
Section~\ref{sec:optimization} was a Redis cache-hit measurement (8.2~ms
mean latency), not a measurement of LLM throughput under concurrent
load} -- despite being reported and discussed in Section~\ref{sec:discussion}
as if it were. The fix appends a fresh UUID nonce to each payload
immediately before firing the concurrent batch, forcing a genuine cache
miss on every one. The corrected measurement is given in
Table~\ref{tab:concurrentfix}.

\begin{table}[htbp]
\caption{Concurrent-Throughput Benchmark: Cache-Hit Bug vs. Corrected}
\label{tab:concurrentfix}
\centering
\begin{tabular}{@{}lll@{}}
\toprule
\textbf{Metric} & \textbf{First-reported (cache-hit bug)} & \textbf{Corrected} \\
\midrule
Concurrent throughput (10 sessions) & 800.3 req/s & \textbf{1.62 req/s} \\
Mean per-request latency & 8.2 ms & \textbf{3.09 s} \\
\bottomrule
\end{tabular}
\end{table}

This correction does not change the sequential-latency, response-diversity,
or blockchain-integrity results in
Sections~\ref{sec:baseline}--\ref{sec:blockchaintestresults}, which were
each measured independently and did not share this code path. It does mean
the Discussion's framing in Section~\ref{sec:discussion} -- that the system
is ``not throughput-starved under concurrent load'' -- is not supported by
this measurement and has been revised accordingly. We report the bug and
the correction explicitly, rather than quietly substituting the corrected
number, because the original figure was published in an earlier revision
of this evaluation and a benchmark bug this large (three orders of
magnitude) is itself a relevant methodological finding for anyone citing or
reproducing this work.

\section{Discussion}
\label{sec:discussion}

The latency/throughput trade-off means this architecture is not a drop-in
replacement for high-volume, low-interaction scanning honeypots -- a
static-path system remains appropriate there. It is instead targeted at
\textbf{high-interaction deception} against a human operator or a
sophisticated, low-volume automated tool, where the multi-second
per-request response latency is plausibly attributable to network
conditions or server load rather than to a honeypot, and where sustained,
contextually consistent engagement yields materially better threat
intelligence than an immediately-abandoned static session.

The distinction between sequential and concurrent throughput (Section
\ref{sec:optimization}, corrected in Section~\ref{sec:realworldrun})
matters for how this trade-off should be read, though the corrected
measurement is less favorable than first reported: a single attacker
session pays the full $\approx$2.06~s mean latency per command against the
real-world dataset (Section~\ref{sec:realworldrun}), and under genuine
concurrent load -- 10 simultaneous sessions, each forced to bypass the
Redis cache -- sustained throughput is \textbf{1.62 requests/second} at a
mean per-request latency of 3.09~s, not the 800~req/s first reported (that
figure was a benchmark-harness bug measuring Redis cache-hit latency; see
Section~\ref{sec:realworldrun}). The corrected number still shows the
architecture serves concurrent sessions without the LLM round-trip time
compounding linearly per session -- 10 concurrent requests complete in
6.16~s wall-clock, not $\approx$20~s ($10 \times 2.06$~s) -- but the
concurrency headroom is real, not the three-orders-of-magnitude figure the
uncorrected benchmark implied. This is the more relevant capacity question
for a real deployment fielding several simultaneous low-volume sessions,
but it does not support a claim of high-volume concurrent capacity.

The blockchain-anchoring design assumes the anchoring chain (or its RPC
endpoint) is itself outside the attacker's reach -- a reasonable assumption
for an internal Ganache/private-chain deployment, less so if a public
chain's RPC endpoint were reachable from the same network segment as a
compromised honeypot host. We did not evaluate anchoring against an
adversary who also compromises the blockchain client, which is out of
scope for this system's threat model (it defends log integrity against
tampering with the \emph{stored log}, not against an attacker who controls
the entire logging pipeline before hashing).

\section{Limitations and Future Work}
\label{sec:limitations}

\begin{itemize}
\item The LLM-judge realism score is an automated proxy; a human-subject
study (as in shelLM) is necessary to validate it against real attacker
perception.
\item Section~\ref{sec:realworldrun} replaced the original 100-payload
hand-curated dataset with a 200-payload dataset built from real,
publicly-documented attack payloads, which is a meaningful improvement over
hand-written strings but is still not a live-traffic capture; it does not
include the full diversity of real-world attacker \emph{behavior}
(multi-step sessions, timing, evasion attempts) a live deployment would
see, only realistic single-payload content.
\item Intent classification is keyword-based and will misclassify
obfuscated or novel attack patterns; a learned classifier is a natural
extension.
\item We did not run a longitudinal live-internet deployment in this
study; results reflect a controlled, synthetic evaluation.
\item The before/after optimization comparison used two different
OpenRouter API keys, because the first key stopped authenticating between
runs for reasons outside this project's control. This does not affect
either run's internal validity, but the two runs were not executed
back-to-back on identical infrastructure state, and the realism-score
change is confounded with a judge-completeness improvement (24/30 scored
vs.\ 30/30) rather than isolating the token-length effect alone -- a
controlled run varying only \texttt{max\_tokens} is needed to attribute
the realism change precisely.
\item The concurrent-throughput measurement used a single burst of 10
simultaneous requests; it was not tested at higher concurrency levels or
under sustained (non-burst) load, so the corrected 1.62~req/s
(Section~\ref{sec:realworldrun}) should be read as a single-burst
demonstration at one concurrency level, not as a validated capacity
ceiling. The original 800~req/s figure this section referenced in earlier
revisions was a benchmark-harness bug (Section~\ref{sec:realworldrun}) and
does not describe the system.
\end{itemize}

\section{Conclusion}

We presented a multi-protocol cyber deception framework that replaces the
static response tables of traditional honeypots with a three-tier,
LLM-backed response engine spanning five distinct decoy services, and that
anchors every log entry's hash on an independent blockchain to make the
resulting forensic trail tamper-evident. Rather than a single unexplained
accuracy figure, we measured and reported the actual trade-off along four
independent axes across three successive evaluation runs: an initial
latency cost of $\approx$4.03~s mean, reduced by a measured
$\approx$3$\times$ via connection pooling and a tuned generation-length cap
to 1.35~s mean (0.4~ms for the static path throughout) against a
hand-curated dataset, rising to 2.06~s mean against a subsequent
200-payload dataset built from real, publicly-documented attack corpora
rather than hand-written strings (Section~\ref{sec:realworldrun}); a
corrected concurrent throughput of 1.62 requests/second at 3.09~s mean
per-request latency under 10 simultaneous sessions, after finding and
fixing a benchmark-harness bug that had silently reported an 800
requests/second Redis cache-hit measurement as if it were LLM throughput;
response diversity of 0.80--1.00 across four of five services against real
payloads (0.70 for the API decoy, a genuine and now documented weakness
against real XXE/Log4Shell payload shapes) versus near-zero for a static
baseline; and a mean LLM-judged realism score stable at 4.00--4.03/5 across
all three runs on complete 30/30-sample scoring -- and we separately
confirmed the blockchain anchoring's tamper-evidence property end-to-end.
We reported the optimization pass's realism cost (a 0.38-point drop,
confounded with a data-completeness improvement), the real-dataset
diversity regression, and the benchmark harness bug alongside the
favorable numbers rather than presenting only the latter. This positions
the system for high-interaction deception against human operators or
low-volume sophisticated tooling -- not as a replacement for the
high-throughput, low-interaction honeypots that remain appropriate for
Internet-scale scanning -- and we have been explicit throughout about
where an LLM-judge score is a proxy rather than a substitute for human
evaluation. Future work should validate the realism findings with human
subjects, replace keyword-based intent classification with a learned
model, and evaluate the system under sustained live-traffic conditions
rather than a synthetic dataset alone.

\begin{thebibliography}{00}

\bibitem{vetterl2018bitter} A. Vetterl and R. Clayton, ``Bitter Harvest: Systematically Fingerprinting Low- and Medium-interaction Honeypots at Internet Scale,'' in \emph{Proc. 12th USENIX Workshop on Offensive Technologies (WOOT '18)}, 2018.

\bibitem{gotcha2023} ``Gotta Catch 'em All: A Multistage Framework for Honeypot Fingerprinting,'' \emph{ACM Digital Threats: Research and Practice}, 2023, arXiv:2109.10652.

\bibitem{honeypotsurvey2026} ``A Review of Honeypots: Fingerprinting Techniques, Detection, and Evasion Mechanisms,'' \emph{MDPI Future Internet}, 2026.

\bibitem{sladic2024shellm} M. Sladi\'{c}, V. Valeros, C. Catania, and S. Garcia, ``LLM in the Shell: Generative Honeypots,'' in \emph{Proc. IEEE EuroS\&PW}, 2024, arXiv:2309.00155.

\bibitem{wang2024honeygpt} Z. Wang et al., ``HoneyGPT: Breaking the Trilemma in Terminal Honeypots with Large Language Model,'' arXiv:2406.01882, 2024.

\bibitem{mckee2023chatbots} F. McKee and D. Noever, ``Chatbots in a Honeypot World,'' arXiv:2301.03771, 2023.

\bibitem{llmhoneypot2024} ``LLM Honeypot: Leveraging Large Language Models as Advanced Interactive Honeypot Systems,'' arXiv:2409.08234, 2024.

\bibitem{blockchainevidence2025} ``Blockchain-Based Evidence Tracking System for Forensic Integrity and Secure Chain of Custody,'' \emph{IEEE}, 2025.

\bibitem{cloudvmforensics2022} ``Blockchain-based tamper-proof and transparent investigation model for cloud VMs,'' \emph{The Journal of Supercomputing}, Springer, 2022, doi:10.1007/s11227-022-04567-4.

\bibitem{peerjthreat2025} ``A novel intelligent cyber threat detection system with blockchain-enforced immutable forensic auditability,'' \emph{PeerJ Computer Science}, 2025.

\bibitem{mitreengage2022} MITRE, ``MITRE Engage: A Framework and Community for Cyber Deception,'' 2022. [Online]. Available: https://engage.mitre.org/

\bibitem{deceptiondecades2021} ``Three Decades of Deception Techniques in Active Cyber Defense -- Retrospect and Outlook,'' \emph{Computers \& Security}, 2021.

\bibitem{bercovitch2011honeygen} M. Bercovitch, M. Renford, L. Hasson, A. Shabtai, L. Rokach, and Y. Elovici, ``HoneyGen: An Automated Honeytokens Generator,'' in \emph{Proc. IEEE International Conference on Intelligence and Security Informatics (ISI)}, 2011.

\bibitem{decoyselection2024} ``A Proactive Decoy Selection Scheme for Cyber Deception using MITRE ATT\&CK,'' arXiv:2404.12783, 2024.

\bibitem{ijnsa2026telemetry} ``Empirical Telemetry-based Metrics for Evaluating Honeypot Realism and Deception Effectiveness,'' \emph{International Journal of Network Security \& Its Applications}, 2026.

\bibitem{cowrie} Cowrie SSH/Telnet Honeypot, GitHub repository. [Online]. Available: https://github.com/cowrie/cowrie

\bibitem{tpot} T-Pot: The All In One Multi Honeypot Platform, GitHub repository. [Online]. Available: https://github.com/telekom-security/tpotce

\bibitem{dionaea} Dionaea Honeypot, GitHub repository. [Online]. Available: https://github.com/DinoTools/dionaea

\bibitem{sokhoneypotllm2025} ``SoK: Honeypots and Large Language Models -- Systematization of a Growing Intersection,'' arXiv:2510.25939, 2025.

\end{thebibliography}

\end{document}
